% ============================================================
% INTERVIEWLEITFADEN V2 — ÜBERARBEITETE FASSUNG
% Revisionsgrundlage: Methodenkritik nach D1–D7 (Witzel 2000;
%   Flick 2011; Flanagan 1954; Przyborski/Wohlrab-Sahr 2021)
%
% Wesentliche Änderungen gegenüber V1:
%   [R1] Block 3 (Führungsrolle/Identität) eliminiert —
%        SDT-externe Konstrukte; Accountability → Block 3 (neu),
%        Identitätserleben emergiert in Block 2a/2b
%   [R2] Block 4 (Rahmenbedingungen) → neu Block 3,
%        auf 12 Min ausgebaut; CIT-Frage ergänzt
%   [R3] S-Fragen in Block 1 von 5 auf 3 reduziert;
%        Präsuppositions-S4 → A-Frage konditionalisiert
%   [R4] Präsuppositionen in 2a-S2, 2c-E, 2c-S1 korrigiert
%   [R5] Block 2b E-Frage: ergebnisverengende Spezifizierung
%        entfernt; S-Fragen von 5 auf 3 reduziert
%   [R6] Einstieg-S-Frage (redundant zur E) gestrichen
%   [R7] Zeitbudget auf realisierbare ~52 Min angepasst
%   [R8] Block 3 (neu): CIT-Frage für Bedingungsanalyse (RQ-B)
%        als S-Frage ergänzt
% ============================================================

% ---- tcolorbox-Stile ---------------------------------------

\tcbset{
  hintbox/.style={
    colback=white,
    colframe=black,
    fontupper=\small\itshape\color{black},
    boxrule=0.6pt,
    leftrule=2.5pt,
    arc=0pt,
    left=6pt, right=6pt,
    top=3pt, bottom=3pt,
    breakable
  }
}

% ---- Hilfsbefehle ------------------------------------------
\newcommand{\Efrage}[1]{%
  \vspace{4pt}%
  \noindent
  \makebox[1.8em][l]{\textcolor{black}{\textbf{E}}}%
  \parbox[t]{\dimexpr\linewidth-1.8em}{%
    \textit{\textbf{#1}}%
  }\par\vspace{4pt}%
}

\newcommand{\Sfrage}[1]{%
  \vspace{2pt}%
  \noindent\hspace{1.8em}%
  \makebox[1.8em][l]{\textcolor{black}{\textbf{S}}}%
  \parbox[t]{\dimexpr\linewidth-3.6em}{#1}\par\vspace{2pt}%
}

\newcommand{\Afrage}[1]{%
  \vspace{2pt}%
  \noindent\hspace{1.8em}%
  \makebox[1.8em][l]{\textcolor{black}{\textbf{A}}}%
  \parbox[t]{\dimexpr\linewidth-3.6em}{#1}\par\vspace{2pt}%
}

\newcommand{\Hinweis}[1]{%
  \begin{tcolorbox}[hintbox]
    $\rightarrow$\enskip #1
  \end{tcolorbox}%
}

\newcommand{\blocksep}{\vspace{0.9em}}

% ============================================================
%  HAUPTDOKUMENT — ab hier einfügen
% ============================================================

\chapter{INTERVIEWLEITFADEN\\\normalsize}

\vspace{0.5em}

% ---- Metadaten-Tabelle -------------------------------------
\begin{tabularx}{\linewidth}{@{}lX@{}}
  \textbf{Projekt}          & \ThesisTitle \\[3pt]
  \textbf{Autor}            & \AuthorName  (Matrikelnummer: \MatriculationNumber) \\[3pt]
  \textbf{Betreuung}        & \Supervisor \\[3pt]
  \textbf{Abgabe}           & \SubmissionDate \\[6pt]
  \textbf{Interviewdatum}   & \underline{\hspace{7cm}} \\[3pt]
  \textbf{Interview-Nummer} & IP-\underline{\hspace{5cm}} \\[3pt]
  \textbf{Dauer (geplant)}  & ca.\ 45–55 Minuten \\[3pt]
  \textbf{Ort / Medium}     & \underline{\hspace{7cm}} \\
\end{tabularx}

\vspace{0.8em}
\noindent\textbf{Legende}

\begin{tabularx}{\linewidth}{@{}llll@{}}
  \textcolor{ecolor}{\textbf{E}} &
    Erzählstimulus / Einstiegsfrage \textit{(kursiv, fett)} \\[4pt]
  \textcolor{scolor}{\textbf{S}} &
    Sondierungsfrage (vertiefend) \\[4pt]
  \textcolor{acolor}{\textbf{A}} &
    Ad-hoc-Frage (situativ, bei Bedarf) \\[4pt]
  $\rightarrow$ &
    Hinweis zur Interviewführung \\
\end{tabularx}

\blocksep

% ============================================================
%  HINWEISE ZUR INTERVIEWFÜHRUNG
% ============================================================
\section{Hinweise zur Interviewführung}

\vspace{0.4em}

\noindent\textbf{Vor dem Interview}

\begin{itemize}[leftmargin=1.5em, itemsep=1pt, topsep=2pt]
  \item[$\rightarrow$] Schriftliche Einverständniserklärung einholen
    (Aufnahme, Datennutzung, Anonymisierung)
  \item[$\rightarrow$] Kurzfragebogen (Anhang) vorab zusenden oder direkt
    zu Beginn ausfüllen lassen
  \item[$\rightarrow$] Voraussetzungen für Aufnahme testen; ruhige,
    störungsfreie Gesprächsatmosphäre sicherstellen
  \item[$\rightarrow$] \textbf{Pilotinterview:} Leitfaden vor Beginn der
    Haupterhebung an 1–2 Personen der Zielgruppe erproben. Prüfen: Werden
    die Fragen wie intendiert verstanden? Öffnen E-Stimuli den Erzählraum?
    Ergebnisse im Postskriptum dokumentieren und Leitfaden ggf.\ anpassen.
    (vgl.\ Witzel 2000)
\end{itemize}

\vspace{0.4em}
\noindent\textbf{Eröffnung des Interviews}

\begin{itemize}[leftmargin=1.5em, itemsep=1pt, topsep=2pt]
  \item[$\rightarrow$] Dank für Teilnahme; kurze Vorstellung des
    Forschungsanliegens (ohne Hypothesen preiszugeben)
  \item[$\rightarrow$] Hinweis: Teilnahme freiwillig, Anonymisierung
    zugesichert, Abbruch jederzeit möglich
  \item[$\rightarrow$] Aufnahmeerlaubnis einholen, Aufnahme starten
\end{itemize}

\vspace{0.4em}
\noindent\textbf{Während des Gesprächs}

\begin{itemize}[leftmargin=1.5em, itemsep=1pt, topsep=2pt]
  \item[$\rightarrow$] Offen und narrativ beginnen, Erzählfluss nicht
    unterbrechen
  \item[$\rightarrow$] S-Fragen erst einsetzen, wenn IP den Erzählimpuls
    nicht selbst aufgreift, nicht als systematische Zweitfrage
    (Witzel 2000, S.~56)
  \item[$\rightarrow$] A-Fragen situativ einsetzen, wenn relevante Themen
    nicht von selbst angesprochen werden
  \item[$\rightarrow$] Aktives Zuhören, Pausen aushalten, kurze
    Rückmeldungen geben (\glqq mhm\grqq{}, \glqq ja\grqq{},
    \glqq interessant\grqq{})
  \item[$\rightarrow$] \textbf{Präsuppositionen im Blick behalten:} Mehrere
    Fragen sind mit \textit{falls\,/\,wenn} konditionalisiert, weil KI-Nutzung
    zwischen IPs stark variieren kann. Diese Konditionalisierungen im
    Gesprächsverlauf beibehalten, nicht glätten.
\end{itemize}

% [R7] Zeitbudget an neue Struktur (4 Blöcke, kein Block 3 alt) angepasst
\vspace{0.4em}
\noindent\textbf{Zeitbudget je Block (Orientierung)}

\vspace{0.2em}
\begin{tabularx}{\linewidth}{@{}lXr@{}}
  \toprule
  \textbf{Abschnitt} & \textbf{Inhalt} & \textbf{Richtwert} \\
  \midrule
  Einstieg  & Rapport, Alltagssituation               & $\approx$\,4 Min \\
  Block 1   & KI-Einsatz, Entscheidungskontext        & $\approx$\,10 Min \\
  Block 2   & Motivationales Erleben (SDT)             & $\approx$\,18 Min \\
  Block 3   & Organisationale Rahmenbedingungen (RQ-B)& $\approx$\,12 Min \\
  Abschluss & Reflexion, Offenes                      & $\approx$\,6 Min \\
  \midrule
  \textbf{Gesamt} &                                   & \textbf{$\approx$\,50 Min} \\
  \bottomrule
\end{tabularx}

\vspace{0.2em}
{\small\itshape Hinweis: Block 3 trägt die Bedingungsanalyse (RQ-B) und ist
  bewusst auf 12~Min ausgebaut. Falls Block~2 länger dauert als geplant,
  Block~2c auf E-Frage + 1~S kürzen, Block~3 darf \textbf{nicht} unterschritten
  werden.}

\vspace{0.4em}
\noindent\textbf{Abschluss und Nachbereitung}

\begin{itemize}[leftmargin=1.5em, itemsep=1pt, topsep=2pt]
  \item[$\rightarrow$] Aufnahme beenden, Dank und Verabschiedung
  \item[$\rightarrow$] Postskriptum direkt nach dem Gespräch anfertigen:
    Besonderheiten, Atmosphäre, nonverbale Signale
  \item[$\rightarrow$] \textbf{Reflexivitätsprotokoll (Pflicht):} Da
    Forscher und Interviewperson demselben professionellen Milieu angehören
    (DACH-Bankensektor, Führungsrolle), besteht erhöhtes Risiko von implizitem
    Vorverständnis und unkritisch geteiltem Deutungswissen. Im Postskriptum
    festhalten: Welche eigenen Annahmen wurden aktiviert? Welche Aussagen der
    IP wirkten überraschend oder wurden vorschnell als selbstverständlich
    übergangen? (vgl.\ Przyborski \& Wohlrab-Sahr 2021, S.~89)
  \item[$\rightarrow$] Transkription nach Dresing \& Pehl (2017),
    Gesprächsführungsprotokoll ablegen
\end{itemize}

\blocksep

% ============================================================
%  GESPRÄCHSEINSTIEG
% ============================================================
\section{Gesprächseinstieg \hfill \normalfont\small\itshape ($\approx$\,4 Min)}

\Efrage{Erzählen Sie mir bitte, wie ein typischer Arbeitstag bei Ihnen als
  Führungskraft aussieht, welche Aufgaben prägen Ihren Alltag am stärksten?}

% [R6] Redundante S-Frage (Zeitanteil) gestrichen — deckungsgleich mit E-Frage
\Afrage{Wenn Sie Ihre Entscheidungsarbeit beschreiben müssten, was würden Sie
  besonders hervorheben?}

\Hinweis{Ziel: Vertrauen aufbauen, Interviewperson situieren, noch keine
  inhaltlichen Vorgaben zu KI machen. S-Frage nur einsetzen, wenn Erzählfluss
  stockt.}

\blocksep

% ============================================================
%  BLOCK 1
% ============================================================
\section{Block 1: Einsatz und Erfahrungen mit generativer KI %
  \hfill \normalfont\small\itshape ($\approx$\,10 Min)}

\Efrage{Ich würde gerne verstehen, welche Rolle generative KI-Tools, wie
  ChatGPT, Microsoft Copilot oder ähnliche, in Ihrer täglichen Führungsarbeit
  spielen. Erzählen Sie mir davon, wie das bei Ihnen konkret aussieht.}

\Sfrage{Welche Tools nutzen Sie, und in welchen Situationen setzen Sie diese ein?}
\Sfrage{Denken Sie an eine Situation, in der die KI-Nutzung besonders gut funktioniert hat, und an eine, in der das eher nicht der Fall war. Was war in diesen Situationen jeweils anders?}
\Afrage{Falls generative KI regelmäßig Teil Ihrer Arbeit ist, wie ist das entstanden? Gab es einen konkreten Ausgangspunkt?}
\Afrage{Gibt es Bereiche Ihrer Führungsarbeit, in denen Sie KI bewusst nicht einsetzen?}

\Hinweis{Offen einsteigen, erst konkrete Nutzungsszenarien erkunden, bevor motivationale Aspekte thematisiert werden.}

\blocksep

\Efrage{Sie haben verschiedene Nutzungssituationen beschrieben. Ich möchte
  auf einen spezifischen Bereich eingehen: Wenn Sie an Entscheidungen denken,
  die Sie vorbereiten müssen, strategische, personelle, operative, wie kommt
  generative KI dabei ins Spiel?}

\Sfrage{Können Sie mir dazu eine konkrete Situation schildern?}
\Sfrage{Was hat sich in diesem Prozess durch die KI-Nutzung verändert –
  wenn überhaupt?}

\Hinweis{Übergang zu Block 2: Fokus gezielt auf Entscheidungsvorbereitungs\-
  prozesse lenken, bevor motivationale Aspekte thematisiert werden.
  Erzählfluss nicht unterbrechen.}

\blocksep

% ============================================================
%  BLOCK 2: MOTIVATIONALE WIRKUNG
% ============================================================
\section{Block 2: Motivationale Wirkung (SDT-Grundbedürfnisse) %
  \hfill \normalfont\small\itshape ($\approx$\,18 Min)}

% ---- 2a) Autonomie -----------------------------------------
\subsubsection{2a) Autonomieerleben \hfill \normalfont\small\itshape ($\approx$\,6 Min)}

\vspace{0.4em}

\Efrage{Wenn Sie an Entscheidungen denken, die Sie mithilfe generativer KI
  vorbereitet haben, wie erleben Sie diesen Prozess für sich selbst?}

\Sfrage{Verändert die KI-Nutzung Ihren eigenen Denkprozess, wie Sie an ein
  Problem herangehen?}
% [R4] Präsupposition ("im Vergleich zu früher") entfernt — konditionaler Vergleich
\Sfrage{Wenn Sie den Abschluss einer KI-gestützten Entscheidung beschreiben –
  wie erleben Sie diesen Moment?}
\Sfrage{Gibt es Situationen, in denen Sie den Prozess als einengend erleben –
  und andere, in denen das nicht der Fall ist? Was unterscheidet diese
  Situationen?}
\Afrage{Wie erleben Sie es, wenn es in Ihrer Organisation konkrete Vorgaben
  zum KI-Einsatz gibt, oder wenn solche fehlen?}

\Hinweis{SDT-Kernbegriff: Autonomie = Erleben von Wahlmöglichkeit,
  Eigenverantwortung und volitionaler Handlungssteuerung. Begriffe
  \textit{Autonomie} und \textit{SDT} werden im Interview \textbf{nicht}
  verwendet.}

\vspace{0.6em}

% ---- 2b) Kompetenz -----------------------------------------
\subsubsection{2b) Kompetenzerleben \hfill \normalfont\small\itshape ($\approx$\,6 Min)}

% [R5] E-Frage: ergebnisverengende Spezifizierung ("in der Art, wie Sie das
%      Ergebnis wahrnehmen") entfernt — Erzählraum für Prozess- UND
%      Ergebniserleben geöffnet
\Efrage{Denken Sie an eine konkrete Arbeitssituation, in der Sie generative
  KI eingesetzt haben. Was hat sich dabei für Sie verändert?}

\Sfrage{Und wie ist das, wenn Sie dieses Ergebnis anderen vorstellen?}
% [R5] S-Fragen von 5 auf 3 reduziert; "besser/schneller" und "fachliches
%      Lernen" zusammengeführt
\Sfrage{Inwiefern verändert die Arbeit mit generativer KI Ihre eigene Wirksamkeit, in dem, was Sie leisten?}
\Sfrage{Gibt es Situationen, in denen KI-Nutzung Ihr Gefühl eigener
  Kompetenz eher untergräbt?}
\Afrage{Wie gehen Sie damit um, wenn ein KI-Vorschlag falsch oder
  unvollständig ist?}

\Hinweis{SDT-Kernbegriff: Kompetenz = Erleben von Wirksamkeit und
  Meisterschaft in herausfordernden Aufgaben. Begriffe nicht im Interview
  verwenden.}

\vspace{0.6em}

% ---- 2c) Soziale Eingebundenheit ---------------------------
\subsubsection{2c) Soziale Eingebundenheit \hfill \normalfont\small\itshape ($\approx$\,6 Min)}

% [R4] Präsupposition ("seitdem KI stärker präsent ist") entfernt;
%      E-Frage für IPs mit sporadischer Nutzung geöffnet
\Efrage{Wie erleben Sie die Zusammenarbeit mit Ihrem Team, wenn generative
  KI im Spiel ist, oder auch, wenn sie es nicht ist?}

% [R4] S-Frage 1 konditionalisiert — kein Voraussetzen von KI-Präsentation
\Sfrage{Wenn Sie im Team auf KI-gestützte Ergebnisse Bezug nehmen –
  wie erleben Sie das?}
\Sfrage{Gibt es Situationen, in denen KI eher verbindend wirkt, und welche,
  in denen sie eher trennend ist?}
\Afrage{Gibt es Situationen, in denen andere unterschiedlich reagieren –
  je nachdem, ob sie wissen, dass KI bei der Erarbeitung beteiligt war?}

\Hinweis{SDT-Kernbegriff: Soziale Eingebundenheit = Gefühl des Dazugehörens,
  bedeutsame Beziehungen zu anderen. Begriffe nicht im Interview verwenden.\\[4pt]
  Übergang zu Block 3: An Konkretes aus Block 2 anknüpfen. Beispiel:
  \glqq Sie haben beschrieben, wie KI Ihren Entscheidungsprozess und Ihre
  Zusammenarbeit beeinflusst. Ich möchte jetzt auf die organisationale Seite
  schauen, welche Rahmenbedingungen das alles prägen.\grqq}

\blocksep

% ============================================================
%  BLOCK 3 (NEU) — ORGANISATIONALE RAHMENBEDINGUNGEN
%  Vormals Block 4; auf 12 Min ausgebaut; RQ-B-Träger
% ============================================================
% [R1][R2] Block 3 alt (Führungsrolle/Identität) eliminiert;
%           ehemaliger Block 4 vorgezogen und erweitert
\section{Block 3: Organisationale Rahmenbedingungen im Bankensektor %
  \hfill \normalfont\small\itshape ($\approx$\,12 Min)}

\Efrage{Wie geht Ihre Organisation mit dem Thema generative KI um, gibt es
  klare Regeln, Empfehlungen oder eher offene Freiräume für den Einsatz?}

\Sfrage{Wie offen ist die Unternehmenskultur gegenüber der Nutzung neuer
  KI-Werkzeuge?}
% [R8] CIT-Frage für Bedingungsanalyse (RQ-B) ergänzt — Flanagan 1954
\Sfrage{Denken Sie an eine Situation, in der der organisationale Rahmen Ihren
  KI-Einsatz gut unterstützt hat, und an eine, in der das nicht der Fall
  war. Was war in diesen Situationen jeweils anders?}
\Sfrage{Wie reagiert Ihr direktes Umfeld, Vorgesetzte, Peers,
  Mitarbeitende, auf Ihre KI-Nutzung?}
\Sfrage{Wie geht Ihre Organisation mit Fehlern um, die durch KI-Outputs
  mitverursacht wurden?}
\Sfrage{Welchen Einfluss hat die regulatorische Umgebung auf Ihre
  Möglichkeiten, KI einzusetzen?}
\Afrage{Gibt es compliance- oder datenschutzbezogene Einschränkungen, die
  Ihre Nutzung konkret begrenzen?}
\Afrage{Was müsste sich verändern, an der Technologie, der Organisation
  oder den Rahmenbedingungen –, damit generative KI Ihre Führungsarbeit
  wirklich gut unterstützt?}
\Afrage{Wie erleben Sie den Umgang mit KI im Vergleich zu anderen
  Instituten, die Sie kennen?}

\Hinweis{Bankensektor-spezifisch: Regulatorik ansprechen, falls nicht von
  der Interviewperson selbst thematisiert. Länderspezifische Referenzen:
  DE: BaFin-Rundschreiben, EBA Guidelines on ML/AI;
  AT: FMA-Leitlinien, DSGVO;
  CH: FINMA-Aufsichtsmitteilungen.
  EU-übergreifend: EU AI Act (Risikoklassifizierung für Kreditinstitute).\\[4pt]
  Die CIT-Frage (S2) ist konstitutiv für die Beantwortung von RQ-B –
  nicht überspringen. (Flanagan 1954)}

\blocksep

% ============================================================
%  GESPRÄCHSABSCHLUSS
% ============================================================
\section{Abschluss und Reflexion \hfill \normalfont\small\itshape ($\approx$\,6 Min)}

% Kurzimpuls Führungsrolle — einziges Überbleibsel des alten Block 3;
% als narrativer Abschlussraum statt eigenem Block
\Efrage{Wenn Sie auf unser gesamtes Gespräch zurückblicken, was würden Sie
  als den Kern Ihres Erlebens mit generativer KI in Ihrer Führungsarbeit
  beschreiben?}

\Sfrage{Inwiefern hat die KI-Nutzung verändert, was Ihre Führungsarbeit für
  Sie im Kern ausmacht, falls das der Fall ist?}

\Afrage{Wenn Sie ein Bild oder einen Begriff wählen müssten, der Ihr
  Verhältnis zu generativer KI treffend einfängt, welcher würde am
  ehesten passen?}
\Afrage{Gibt es etwas, das wir noch nicht angesprochen haben, das Ihnen aber
  wichtig erscheint?}

\Hinweis{Aufnahme beenden. Dank und Verabschiedung. Hinweis auf
  Anonymisierung und Weiterverwendung der Daten. Postskriptum \textbf{und}
  Reflexivitätsprotokoll unmittelbar nach dem Gespräch anfertigen:
  Besonderheiten, nonverbale Signale, Gesprächsatmosphäre, eigene aktivierte
  Vorannahmen.}

\blocksep

% ============================================================
%  ANHANG: KURZFRAGEBOGEN
% ============================================================
\section*{Anhang: Kurzfragebogen (Soziodemografika)}
\textit{Bitte vor oder unmittelbar nach dem Interview ausfüllen. Alle Angaben
  sind anonym und dienen ausschließlich der Beschreibung der Stichprobe.}

\noindent\textbf{Angaben zur Person und Tätigkeit}

\vspace{0.4em}

\begin{tabularx}{\linewidth}{@{}lX@{}}
  \toprule
  \textbf{Interview-Nummer}
    & IP-\underline{\hspace{7cm}} \\[5pt]
  \textbf{Aktuelle Position / Bezeichnung}
    & \underline{\hspace{7cm}} \\[5pt]
  \textbf{Hierarchieebene} & $\square$ Teamleitung \quad
      $\square$ Gruppenleitung \\ &
      $\square$ Abteilungsleitung \quad
      $\square$ Bereichsleitung \\[5pt]
  \textbf{Dauer in aktueller Führungsrolle}
    & \underline{\hspace{4cm}} Jahre \\[5pt]
  \textbf{Führungsspanne (direkt unterst.)}
    & \underline{\hspace{4cm}} Mitarbeitende \\
  \bottomrule
\end{tabularx}

\vspace{0.8em}

\noindent\textbf{Angaben zur Organisation}

\vspace{0.4em}

\begin{tabularx}{\linewidth}{@{}lX@{}}
  \toprule
  \textbf{Typ des Kreditinstituts} &
    $\square$ Universalbank \quad
    $\square$ Privatbank \\ &
    $\square$ Sparkasse \quad
    $\square$ Genossenschaftsbank \\ &
    $\square$ Direktbank \quad
    $\square$ Sonstige \\[5pt]
  \textbf{Land} &
    $\square$ Deutschland \quad
    $\square$ Österreich \quad
    $\square$ Schweiz \\[5pt]
  \textbf{Institutsgröße (MA gesamt, ca.)} &
    $\square$ ${<}$\,500 \quad
    $\square$ 500–5.000 \\ &
    $\square$ 5.001–20.000 \quad
    $\square$ ${>}$\,20.000 \\
  \bottomrule
\end{tabularx}

\vspace{0.8em}

\noindent\textbf{Angaben zur KI-Nutzung}

\vspace{0.4em}

\begin{tabularx}{\linewidth}{@{}lX@{}}
  \toprule
  \textbf{Nutzungshäufigkeit KI-Tools} &
    $\square$ täglich \quad
    $\square$ mehrmals/Woche \\ &
    $\square$ wöchentlich \quad
    $\square$ selten \\ &
    $\square$ gar nicht \\[5pt]
  \textbf{Genutzte Tools (offene Nennung)} &
    \underline{\hspace{8cm}} \\[5pt]
  \textbf{KI-Nutzung dienstlich erlaubt?} &
    $\square$ Ja, offiziell eingeführt \\ &
    $\square$ Ja, geduldet / informell \\ &
    $\square$ Nein / unklar \\
  \bottomrule
\end{tabularx}

\vspace{0.8em}

\begin{tabularx}{\linewidth}{@{}lX@{}}
  \toprule
  \textbf{Anmerkungen / Ergänzungen} &
    \underline{\hspace{8cm}} \\[5pt]
                                     &
    \underline{\hspace{8cm}} \\
  \bottomrule
\end{tabularx}

\vspace{1em}